% \iffalse meta-comment
%
% RWTH-CI -- Corporate Design for RWTH Aachen University
% ----------------------------------------------------------------------------
%
%  Copyright (C) 2025–2026 by Marei Peischl <rwth-ci@peitex.de> and Lukas C. Bossert
%
% ============================================================================
% This work may be distributed and/or modified under the
% conditions of the LaTeX Project Public License, either version 1.3c
% of this license or (at your option) any later version.
% The latest version of this license is in
% http://www.latex-project.org/lppl.txt
% and version 1.3c or later is part of all distributions of LaTeX
% version 2008/05/04 or later.
%
% This work has the LPPL maintenance status `maintained'.
%
% The Current Maintainer of this work is
%   Marei Peischl rwth-ci@peitex.de
%
% The development repository can be found at
% https://gitlab.git.nrw/rwth-it-center/rwth-latex-templates/rwth-ci
% Please use the issue tracker for feedback!
%
% ============================================================================
%
% \fi
% \iffalse
%<*driver>
\ProvidesFile{rwth-colors.dtx}
  [2026-04-03 v1.2 Color setup for RWTH-CI, Corporate Design of RWTH Aachen University]
%</driver>
%<@@=ptxcd>
%<*identification>
\NeedsTeXFormat{LaTeX2e}[2022-10-01]
%<package>\ProvidesExplPackage{rwth-colors}{2026-04-03}{1.2}{Color setup for RWTH-CI}
%</identification>
%<*driver>
\providecommand*{\PtxcdDocDTXfiles}{rwth-colors.dtx}
\providecommand*{\PtxcdDocTitle}{\pkg{rwth-colors}~--~\LaTeX{}~color~setup~for~RWTH-CI}
% \iffalse meta-comment
%
% RWTH-CI -- Corporate Design for RWTH Aachen University
% ----------------------------------------------------------------------------
%
%  Copyright (C) 2025–2026 by Marei Peischl <rwth-ci@peitex.de> and Lukas C. Bossert
%
% ============================================================================
% This work may be distributed and/or modified under the
% conditions of the LaTeX Project Public License, either version 1.3c
% of this license or (at your option) any later version.
% The latest version of this license is in
% http://www.latex-project.org/lppl.txt
% and version 1.3c or later is part of all distributions of LaTeX
% version 2008/05/04 or later.
%
% This work has the LPPL maintenance status `maintained'.
%
% The Current Maintainer of this work is
%   Marei Peischl <rwth-ci@peitex.de>
%
% The development repository can be found at
% https://gitlab.git.nrw/rwth-it-center/rwth-latex-templates/rwth-ci
% Please use the issue tracker for feedback!
% ============================================================================
%
% \fi
%
% \iffalse
%<*driver>
\ProvidesFile{rwth-ci.dtx}
  [2026-04-03 v1.2 Fundamental configuration for RWTH-CI]
%</driver>
%<@@=ptxcd>
%<*identification>
%<package|class>\NeedsTeXFormat{LaTeX2e}[2022/10/01]
%<package>\ProvidesExplPackage{rwth-layout}{2026-04-03}{1.2}{Fundamental configuration for for RWTH-CI}
%</identification>
%<*driver>
\documentclass[lm-default=false,cs-break-nohyphen]{l3doc}

\usepackage[highlightmode=immediate]{minted}
\usepackage{accsupp}
\setminted{
	tabsize=2,
	breaksymbol={\BeginAccSupp{method=escape,ActualText={}}\tiny\ensuremath{\hookrightarrow}\EndAccSupp{}}
}
\newenvironment{examplecode}{%
		\VerbatimEnvironment%
		\fvset{gobble=0}%
		\begin{VerbatimOut}[gobble=1]{examplecode.out}%
}{%
		\end{VerbatimOut}%
		\inputminted[autogobble, escapeinside=||]{latex}{examplecode.out}%
}

\ExplSyntaxOn
\makeatletter
\tl_new:N \l__ptxtools_doc_values_tl
\newcommand*{\codefamily}{\MacroFont}
\DeclareTextFontCommand{\codefont}{\codefamily}
\providecommand{\option}[1]{\codefont{#1}}
\newcommand{\optionvalue}[1]{\codefont{#1}}

\cs_new:Npn \__ptxctools_parse_key_option:w #1 #2 #3 #4 = #5 \q_stop {
  \begingroup
  \let\PrintDescribeOption\PrintDescribeKeyOption
  \tl_set:Nn \l__ptxtools_doc_values_tl {\small\IfBooleanTF{#3}{\textsf{#5}}{\textsf{(\codefont{#5})}}}
  \noindent\DescribeOption{#4}\hfill
  \small
  \tl_if_blank:nTF {#2} {\meta{\textrm{initially~unset}}} {
    \IfBooleanTF{#1}
    {#2}
    {(default:~\codefont{#2})}
  }\strut\par
  \endgroup
}

\newcommand*{\PrintDescribeKeyOption}[1]{
  {\MacroFont #1=}\rlap{\hskip\marginparsep\l__ptxtools_doc_values_tl}
}

\bool_new:N \l_ptxtools_last_was_describe_bool
\NewDocumentCommand{\DescribeKeyOption}{smms}{
  \bool_if:NTF \l_ptxtools_last_was_describe_bool {
    \vskip -\lastskip
  }{
    \par
    \dim_compare:nNnF {\parskip} > {\c_zero_dim} {\medskip}
  }
  \__ptxctools_parse_key_option:w #1 {#3} #4 #2 \q_stop

  \bool_set_true:N \l_ptxtools_last_was_describe_bool
  \AddToHookNext{para/begin}{
    \bool_set_false:N \l_ptxtools_last_was_describe_bool
    \OmitIndent
  }
  \smallskip
}

% to support small layout adjustments for the texdoc version
\newcommand*{\IfOnlyDocumentation}[1]{
  \bool_if_exist:NT \g__codedoc_typeset_implementation_bool {
    \bool_if:NF \g__codedoc_typeset_implementation_bool {#1}
  }
}

\NewDocumentCommand{\PtxtoolsLoadDocDTX}{m}{
  \exp_args:No \clist_map_inline:nn {#1} {\DocInput{##1}}
}

\ExplSyntaxOff

\usepackage{underscore-ltx}

\DeclareRobustCommand{\issue}[1]{%
  {\DeclareFieldFormat{citeurlpostfix}{\href{##1/issues/#1}{\##1}}\citefield{othr-ci-repo}[citeurlpostfix]{url}}%
}
\newcommand*{\version}[1]{v#1}

\usepackage{tabularx}
\usepackage{biblatex}
\addbibresource{DEMO-RWTHBibliography.bib}

\usepackage{pifont}% Zapf-Dingbats Symbols
\newcommand*{\FeatureTrue}{\ding{52}}
\newcommand*{\FeatureFalse}{\ding{56}}

% Load RWTH-CI packages
\RequirePackage{rwth-colors}
\EnableCrossrefs

\CodelineIndex
\RecordChanges

\providecommand*{\PtxcdDocTitle}{\pkg{\jobname} -- \LaTeX{} using the Corporate Design of RWTH Aachen University}

\GetFileInfo{\jobname.dtx}
\DoNotIndex{\newcommand,\newenvironment}
\title{\PtxcdDocTitle^^A
  \thanks{This file describes version \fileversion, last revised \filedate.}
}
\author{Marei Peischl\thanks{Email: \href{mailto:rwth-ci@peitex.de}{rwth-ci@peitex.de}}}
\date{\fileversion~of \filedate}

\renewcommand*{\partname}{Part}

% avoid conflict with | in bibliography
\RemoveFromHook{begindocument}[l3doc]
\AddToHook{begindocument/end}{
	\MakeShortVerb \"
	\MakeShortVerb \|
}

\begin{document}

\maketitle
\OnlyDescription
\DisableImplementation

\tableofcontents

\clearpage

\expandafter\ifx\csname PtxcdDocDTXfiles\endcsname\relax
	\providecommand*{\PtxcdDocDTXfiles}{rwth-ci.dtx,rwth-fonts.dtx,rwth-colors.dtx,rwth-letter.dtx,rwth-beamer.dtx,rwth-sciposter.dtx}
	\part{Global configuration and auxiliary packages}
	\PtxtoolsLoadDocDTX{rwth-ci.dtx,rwth-fonts.dtx,rwth-colors.dtx}

	\clearpage
	\part{Document Types}
	\PtxtoolsLoadDocDTX{rwth-letter.dtx,rwth-beamer.dtx,rwth-sciposter.dtx}
\else
	\PtxtoolsLoadDocDTX{\PtxcdDocDTXfiles}
\fi

% Don't include Implementation on CTAN
\expandafter\ifx\csname DoNotTypesetImplementation\endcsname\relax
  \DisableDocumentation\EnableImplementation
  \begin{implementation}
  \part{Implementation}
  \end{implementation}
  \PtxtoolsLoadDocDTX{\PtxcdDocDTXfiles}
\fi

\PrintChanges
\printbibliography
% Restore \code to be a wrapper for index numbers
\begingroup
\renewcommand*\code[1]{#1}
\PrintIndex
\endgroup

\end{document}
% \fi
% \begin{documentation}
%
% \end{documentation}
% \iffalse
%</driver>
% \fi
%    \begin{macrocode}
%<class|package|(init&logo)>\tl_new:N \g_ptxcd_logofile_tl
%    \end{macrocode}
% \iffalse
%<*init&department>
% \fi
%    \begin{macrocode}
\tl_if_exist:NF \g_@@_config_prefix_tl {
	\tl_new:N \g_@@_config_prefix_tl
	\tl_gset:Nn \g_@@_config_prefix_tl {rwth-}
}
%    \end{macrocode}
% \iffalse
%</init&department>
%<*class|package>
% \fi
% \begin{implementation}
%    \begin{macrocode}
%<!beamer>\prop_new:N \g_ptxcd_clsopts_prop
%<!beamer>\int_new:N \g_ptxcd_paper_int
%<sciposter>\bool_new:N \g_@@_logo_bottom_bool
\str_const:Nn \c_@@_base_str {
%<beamer>  beamer
%<sciposter>  sciposter
%<package>  layout
}
%    \end{macrocode}
% \iffalse
%</class|package>
%<*class|(init&department)>
% \fi
% \changes{v1.2-dev}{2026-02-17}{Fix department config for letters and prepare to support additional packages.}
%    \begin{macrocode}
\msg_if_exist:nnF {rwth-ci} {no-departmentconfig-found} {
	\msg_new:nnn {rwth-ci} {no-departmentconfig-found} {
		The~package~has~been~configured~to~load~the~config~file~“#1”.\\
		I~can't~find~this~file.~Skipping.
	}
}
\hook_gput_code:nnn {
%<class>	class/\@currname/after
%<!class>	package/\@currname/after
}{
  \@currname/department-config
}{
  \file_if_exist_input:nF {\g_@@_config_prefix_tl\g_ptxcd_department_str.cfg} {
    \str_if_eq:eeF {\g_@@_config_prefix_tl\g_ptxcd_department_str} {rwth-default} {
      \msg_note:nne {rwth-ci} {no-departmentconfig-found} {\g_@@_config_prefix_tl\g_ptxcd_department_str.cfg}
    }
  }
}
%    \end{macrocode}
% \iffalse
%</class|(init&department)>
%<*class|package>
% \fi
%    \begin{macrocode}
\exp_args:Ne \keys_define:nn {ptxcd/\c_@@_base_str} {
%    \end{macrocode}
% \iffalse
%</class|package>
%<*options>
% \fi
% \end{implementation}
% \begin{documentation}
% \section{General configuration options}
% \DescribeKeyOption{logofile=\meta{file name/path}}{rwth_mx}
% Is used to configure specific logos.
% The logos are reserved to be used by RWTH departments and may therefore not be published within this bundle.
%
% Depending on the color mode (see \autoref{sec:colormode}) the actual file searched for will be extended by the suffices |_cmyk| or |_rgb|.
% Within this document the file |rwth_mx_cmyk| was used.
% \end{documentation}
% \begin{implementation}
% \begin{optionenv}{logofile}
% \changes{v1.0}{2025-12-18}{init .pdf extension if no extension was found}
%    \begin{macrocode}
  logofile .code:n = {
    \file_parse_full_name:nNNN {#1} \g_tmpa_tl \l_tmpa_tl \l_tmpb_tl
    \tl_gset:Ne \g_ptxcd_logofile_tl {\tl_if_empty:NF \g_tmpa_tl  {\g_tmpa_tl/}\l_tmpa_tl}
    \tl_if_empty:NTF \l_tmpb_tl {
      \tl_if_empty:NT \g_@@_logofile_ext_tl {
        \tl_gset:Nn \g_@@_logofile_ext_tl {.pdf}
      }
    }{
      \tl_gset_eq:NN \g_@@_logofile_ext_tl \l_tmpb_tl
    }
  },
%<class>  logofile .initial:n = ,
%    \end{macrocode}
% \end{optionenv}
% \end{implementation}
% \begin{documentation}
% \DescribeKeyOption{logofile-ext=\meta{extension}}{.pdf}
% \end{documentation}
% \begin{implementation}
% \begin{optionenv}{logofile-ext}
%    \begin{macrocode}
  logofile-ext .tl_gset:N = \g_@@_logofile_ext_tl,
%<class>  logofile-ext .initial:n = .pdf,
%    \end{macrocode}
% \end{optionenv}
% \end{implementation}
% \begin{documentation}
% \DescribeKeyOption{accept-missing-logos=\meta{boolean}}{false but switched to true, if no logofile is selected}
% As the logo files sadly may not be published within this bundle, the default setting allows all files to be compiled without a logo.
% There will be a warning and the logo postion will be marked showing the desired filename.
% \end{documentation}
% \begin{implementation}
% \begin{optionenv}{accept-missing-logos}
%    \begin{macrocode}
  accept-missing-logos .bool_gset:N = \g_@@_logo_workaround_bool,
%<class>  accept-missing-logos .initial:n = false,
  accept-missing-logos .usage:n = load,
  accept-missing-logos .default:n = true,
%    \end{macrocode}
% \end{optionenv}
% \end{implementation}
% \iffalse
%</options>
%<*options|fonts>
% \fi
% \begin{implementation}
% undocumented feature
% \begin{optionenv}{LaTeX-logo}
%    \begin{macrocode}
  LaTeX-logo .bool_gset:N = \g_@@_use_LaTeX_Logo_bool,
%<fonts>  LaTeX-logo .initial:n = false,
%<fonts>  LaTeX-logo .usage:n = load,
%    \end{macrocode}
% \end{optionenv}
% \end{implementation}
% \iffalse
%</options|fonts>
%<*options>
% \fi
% \begin{documentation}
% \DescribeKeyOption{department=\meta{department}}{default}
% RWTH-CI supports department specific configuration files.
% The \optionvalue{default} is used to detect if a user specific setting was provided.
% \DescribeKeyOption{departmentconfigprefix=\meta{prefix}}{rwth-}
% By default this setting will search for files called |rwth-|\meta{department}|.cfg|, where |rwth-| can be changed using \option{departmentconfigprefix}.
% This option could also be used to load config files from  specific subdirectories.
% This might be useful if a specific department wants some internal templates to be shared separately.
% \end{documentation}
% \begin{implementation}
% \begin{optionenv}{department}
%    \begin{macrocode}
  department .choice:,
  department / default .code:n = \str_gset:Nn \g_ptxcd_department_str {default},
%<class>  department .initial:n = default,
  department / unknown .code:n = {
      \str_gset:NV \g_ptxcd_department_str \l_keys_value_tl
    },
%    \end{macrocode}
% \end{optionenv}
% \begin{optionenv}{departmentconfigprefix}
%    \begin{macrocode}
  departmentconfigprefix .tl_gset:N = \g_@@_config_prefix_tl,
%<class>  departmentconfigprefix .initial:n = rwth-,
%    \end{macrocode}
% \end{optionenv}
% \end{implementation}
% \begin{documentation}
% \DescribeKeyOption{colormode=RGB/cmyk}{depens on document type}
% Selects the color scheme to be converted to.
% This will also affect the logo selection.
% \end{documentation}
% \begin{implementation}
% \begin{optionenv}{colormode}
%    \begin{macrocode}
  colormode .code:n = \PassOptionsToPackage{colormode=#1}{rwth-colors},
  RGB .code:n = \PassOptionsToPackage{colormode=RGB}{rwth-colors},
  cmyk .code:n = \PassOptionsToPackage{colormode=RGB}{rwth-colors},
%    \end{macrocode}
% \end{optionenv}
%    \begin{macrocode}
%<!titlegraphic-opt&(init-department|class)>  titlegraphic .code:n = \PassOptionsToPackage{titlegraphic={#1}}{beamerinnerthemeRWTH},
%    \end{macrocode}
% \end{implementation}
% \iffalse
%<*sciposter>
% \fi
% \begin{documentation}
% \DescribeKeyOption{logo=top/bottom}{top}
% Allows to change the logo position. This option is currently only supported by the \cls{rwth-sciposter}.
% \end{documentation}
% \begin{implementation}
% \begin{optionenv}{logo}
%    \begin{macrocode}
  logo .choice:,
  logo / bottom .code:n = { \bool_gset_true:N \g_@@_logo_bottom_bool },
  logo / top .code:n = { \bool_gset_false:N \g_@@_logo_bottom_bool },
  logo / head .meta:n = {logo=top},
  logo / foot .meta:n = {logo=bottom},
  logo .initial:n = {top},
%    \end{macrocode}
% \end{optionenv}
% \end{implementation}
% \iffalse
%</sciposter>
%<*letter|sciposter|package>
% \fi
% \begin{documentation}
% \DescribeKeyOption{raggedright=\meta{boolean}}{true}
% The CI requires all text to be set ragged right.
% This option can be used to disable this setting to avoid conflicts with additional packages or configuration to adjust the placement of elements.
% In general this setting will be enabled at |\begin{document}| but for \pkg{rwth-letter} the option \option{letterhooks} allows to limit this setting for the letter environment.
% \end{documentation}
% \begin{implementation}
% \begin{optionenv}{raggedright}
%    \begin{macrocode}
  raggedright .bool_gset:N = \g_@@_raggedright_bool,
  raggedright .initial:n = true,
  raggedright .default:n = true,
  raggedright .usage:n = load,
%    \end{macrocode}
% \end{optionenv}
% \iffalse
%</letter|sciposter|package>
%</options>
%<*titlegraphic-opt>
% \fi
% \end{implementation}
% \begin{documentation}
% \DescribeKeyOption{titlegraphic=1:3,2:3,false}{1:3}
% The titlegraphic can be used in different sizes.
% This option is currently only supported by \cls{rwth-beamer}.
%
% This option can be changed locally.
% Therefore the |\titlepage| command was extended to accept an optional argument.
% \begin{examplecode}
% \titlepage[titlegraphic=1:3]
% \end{examplecode}
% \end{documentation}
% \begin{implementation}
% \begin{optionenv}{titlegraphic}
%    \begin{macrocode}
  titlegraphic-ratio .str_set:N = \l__ptxcd_titlegraphic_ratio_str,
  titlegraphic .multichoice:,
  titlegraphic / 1:3 .meta:n ={titlegraphic-ratio={1:3}},
  titlegraphic / 2:3  .meta:n = {titlegraphic-ratio={2:3}},
  titlegraphic / false .meta:n = {titlepage=plain},
  titlegraphic / true .meta:n = {titlepage=graphics},
  titlegraphic .initial:n = 1:3,
%    \end{macrocode}
% \end{optionenv}
% \end{implementation}
% \iffalse
%</titlegraphic-opt>
%<*options>
%<*titlesepline>
% \fi
% \begin{documentation}
% \DescribeKeyOption{titlesepline=\meta{boolean}}{false}
% This option is currently only supported by \cls{rwth-beamer}.
% \end{documentation}
% \begin{implementation}
% \begin{optionenv}{titlesepline}
%    \begin{macrocode},
  titlesepline .bool_gset:N = \g_@@_titlesepline_bool,
  titlesepline .initial:n = true,
  titlesepline .default:n = true,
%    \end{macrocode}
% \end{optionenv}
% \end{implementation}
% \iffalse
%</titlesepline>
%<*print>
% \fi
% \begin{implementation}
% fontsize
% \begin{optionenv}{fontsize}
%    \begin{macrocode}
  fontsize .code:n = \prop_gput:Nnn \g_ptxcd_clsopts_prop {fontsize} {#1},
  fontsize .initial:n =
%<beamer> 9 pt,% beamer
%<sciposter> 23bp,
%    \end{macrocode}
% \end{optionenv}
% \end{implementation}
% \begin{documentation}
% \iffalse\DescribeKeyOption{paper=a\meta{n}/portrait/landscape}{\meta{depends on document type}}\fi
% \DescribeKeyOption{paper=portrait/landscape}{\meta{depends on document type}}
% The supported paper sizes as well as the default setting depends on the document type.
% \iffalse Standard print documents use \optionvalue{a4} and support the fulle range from \optionvalue{a0} to \optionvalue{a6}. \fi
% The default values for posters is \optionvalue{a0} and \optionvalue{portrait}.
% \end{documentation}
% \begin{implementation}
% \changes{v1.1}{2026-04-03}{Add support for landscape and portrait als boolean options.}
% \begin{optionenv}{paper}
%    \begin{macrocode}
  paper .multichoices:nn = {
%<sciposter>      a0%%,a1,a2,a3
%<poster>      a0,a1,a2,a3,a4
%<!(sciposter|poster)>      a0,a1,a2,a3,a4,a5,a6
}{
  \int_gset_eq:NN \g_ptxcd_paper_int  \l_keys_choice_int
%<!(sciposter|poster)>      \exp_args:Nx \PassOptionsToPackage{paper=\l_keys_choice_tl}{typearea}
      \exp_args:Nx \PassOptionsToPackage{\l_keys_choice_tl paper}{geometry}
    },
  paper .initial:n =
%<!(sciposter|poster)>  a4,
%<sciposter|poster>  a0,
%    \end{macrocode}
% \iffalse
%<*sciposter>
% \fi
%    \begin{macrocode}
  paper / landscape .code:n = \PassOptionsToPackage{landscape}{geometry},
  paper / portrait .code:n = \PassOptionsToPackage{portrait}{geometry},
  landscape / .choice:,
  landscape / true .meta:n = {paper = landscape},
  landscape / false .meta:n = {paper = portrait},
  portrait / .choice:,
  portrait / true .meta:n = {paper = portrait},
  portrait / false .meta:n = {paper = landscape},
%    \end{macrocode}
% \end{optionenv}
% \iffalse
%</sciposter>
%</print>
%<*class&!beamer>
% \fi
%    \begin{macrocode}
  unknown .code:n = {\prop_gput:NVn \g_ptxcd_clsopts_prop \l_keys_key_tl {#1}},
%    \end{macrocode}
% \iffalse
%</class&!beamer>
%</options>
% \fi
%    \begin{macrocode}
%<processoptions|package>}
%    \end{macrocode}
% \iffalse
%<*processoptions|init-department>
% \fi
%    \begin{macrocode}
\str_if_exist:NF \g_ptxcd_department_str {
  \str_new:N \g_ptxcd_department_str
  \str_gset:Nn \g_ptxcd_department_str {default}
}
%    \end{macrocode}
% \iffalse
%</processoptions|init-department>
%<*!theme>
%<*processoptions|(beamer&buildOptions)>
% \fi
% Allow to add options via l3build |typesetcmds|.
% Be aware that in this case also the group has to be expanded.
%    \begin{macrocode}
\cs_if_exist:NT \ptxcdDefaultBuildOptions {
  \exp_args:Ne \keys_set:ne { ptxcd /
%<class|beamer>    \c_@@_base_str
%<letter>    letter
  } {
    \ptxcdDefaultBuildOptions
  }
}
%    \end{macrocode}
% \iffalse
%</processoptions|(beamer&buildOptions)>
%</!theme>
%<*processoptions|package>
% \fi
%    \begin{macrocode}
%<class|package>\ProcessKeyOptions[ptxcd/\c_@@_base_str]
%    \end{macrocode}
% \iffalse
%</processoptions|package>
%<ExplSyntax>\ExplSyntaxOn
%<*logoFallback>
% \fi
% \changes{v0.2}{2025-11-12}{enable accept-missing-logos for the default logo setup}
%    \begin{macrocode}
\tl_if_empty:NT \g_ptxcd_logofile_tl {
  \msg_new:nnn {rwth-ci} {no-logo-selected} {You~did~not~select~a~logofile.\\
    I~will~try~to~use~the~example~logo~`rwth_mx`~and~enable~the~`accept-missing-logos`}
  \msg_warning:nn {rwth-ci} {no-logo-selected}
  \tl_gset:Nn \g_ptxcd_logofile_tl {rwth_mx}
  \tl_gset:Nn \g_@@_logofile_ext_tl {.pdf}
  \bool_gset_true:N \g_@@_logo_workaround_bool
}
%    \end{macrocode}
% Define a macro to insert the logo workaround to only need this definition once.
% \changes{v0.2}{2025-11-12}{Add warning if logo fallback is active.}
%    \begin{macrocode}
\msg_new:nnn {rwth-ci} {missing-logofile} {%
  The~logofile~(#1)~could~not~be~found.\\
  Please~check~if~if~can~be~found~by~the~compiler.
}
\cs_new:Nn \_@@_insert_logo:n {
%    \end{macrocode}
% TODO: Check if rgb is a good default for logo selection.
% Try to detect if colormode specific file is available.
%    \begin{macrocode}
  \file_if_exist:nTF {\g_ptxcd_logofile_tl\tl_if_empty:NTF \g_@@_colormode_tl {_rgb} {_\g_@@_colormode_tl}\g_@@_logofile_ext_tl } {
    \bool_set_true:N \l_tmpa_bool
    \bool_set_true:N \l_tmpb_bool
  }{
    \bool_set_false:N \l_tmpa_bool
    \file_if_exist:nTF {\g_ptxcd_logofile_tl\g_@@_logofile_ext_tl}
      {\bool_set_true:N \l_tmpb_bool}
      {\bool_set_false:N \l_tmpb_bool}
  }
  \bool_if:NTF \g_@@_logo_workaround_bool {
    \bool_if:NTF \l_tmpb_bool
    \use_i:nn
    \use_ii:nn
  }{
    \use_i:nn
  }
  {\includegraphics[height=#1]{\g_ptxcd_logofile_tl\bool_if:NT \l_tmpa_bool {\tl_if_empty:NTF \g_@@_colormode_tl {_rgb} {_\g_@@_colormode_tl}}\g_@@_logofile_ext_tl}}
  {
%<!beamer>    \frame
%<beamer>    \framelatex
    {\parbox[c][#1][c]{2.186\dimexpr#1}{\centering<Missing~Logo~file~\tl_to_str:N \g_ptxcd_logofile_tl>}}
    \msg_warning:nnV {rwth-ci} {missing-logofile} \g_ptxcd_logofile_tl
  }
}
%    \end{macrocode}
% \iffalse
%</logoFallback>
%<*titlegraphic-setup>
% \fi
% \end{implementation}
% \begin{documentation}
% \section{RWTH-CI specific title elements}
% All templates except the \pkg{rwth-letter}, are configured to use |\title|, |author|, \dots for configuration of the title block/page.
% There are some additional elements which can be used depending on the document type.
%
% \DescribeMacro{\titlegraphic}
% Accepts any \LaTeX{} content, does not have to be an image.
% This is placed flush with the top corner in the main part of the title page.
% This macro is usually used to place a graphic:
% \begin{examplecode}
%   \titlegraphic{\includegraphics[width=\width]{example-image-plain}}
% \end{examplecode}
% \DescribeMacro{\width}\DescribeMacro{\height}
% \cs[no-index]{width} and \cs[no-index]{height} can be used to select the appropriate size of the image.
% There also is a starred variant \cs[no-index]{titlegraphic*}, to allow automatic scaling and any necessary clipping to fill the reserved space. This is done automatically using \pkg{trimclip}.
% \begin{examplecode}
%   \titlegraphic*{\includegraphics{example-image-plain}}
% \end{examplecode}
% \end{documentation}
% \begin{implementation}
%    \begin{macrocode}
\RequirePackage{trimclip}
\box_new:N \l_@@_titlegraphic_box
%<!beamer>\NewDocumentCommand
%<beamer>\RenewDocumentCommand
{\titlegraphic}{sm}{
  \IfBooleanTF{#1}{
%<!beamer>    \tl_gset:Nn  \g_ptxcd_titleimage_code_tl
%<beamer>    \def\inserttitlegraphic
    {
      \hbox_set:Nn \l_@@_titlegraphic_box {\raisebox{\depth}{#2}}
      \box_resize_to_wd:Nn \l_@@_titlegraphic_box {\width}
      \dim_compare:nTF {\box_ht:N \l_@@_titlegraphic_box -\height> \c_zero_dim}
      {
        \dim_set:Nn \l_tmpa_dim {.5\box_ht:N \l_@@_titlegraphic_box - .5\height}
        \clipbox{0pt~\dim_eval:n{\l_tmpa_dim}~0pt~\dim_eval:n{\l_tmpa_dim}}{\box_use:N \l_@@_titlegraphic_box}
      }{
        \box_resize_to_ht:Nn \l_@@_titlegraphic_box {\height}
        \dim_set:Nn \l_tmpa_dim {(\box_wd:N \l_@@_titlegraphic_box - \width) / 2}
        \clipbox{\dim_eval:n{\l_tmpa_dim}~0pt~\dim_eval:n{\l_tmpa_dim}~0pt}{\box_use:N \l_@@_titlegraphic_box}
      }
    }
  }{
%<!beamer>    \tl_gset:Nn  \g_ptxcd_titleimage_code_tl
%<beamer>    \def\inserttitlegraphic
    {#2}
  }
}
%    \end{macrocode}
% \iffalse
%</titlegraphic-setup>
% \fi
%    \begin{macrocode}
%<ExplSyntax>\ExplSyntaxOff
%    \end{macrocode}
% \iffalse
%<*package>
% \fi
%    \begin{macrocode}
\PassOptionsToPackage{fontdefaults=true}{rwth-fonts}
\RequirePackage{rwth-fonts}
\RequirePackage{rwth-colors}
\bool_if:NT  \g_@@_raggedright_bool {
  \hook_gput_code:nnn {begindocument/end} {rwth-ci/raggedright} {\raggedright}
}
\RequirePackage{scrlayer-scrpage}
\RequirePackage{graphicx}
\DeclareNewLayer[background,align=tl,voffset=0mm, hoffset=20mm,
  mode=picture,
  width=\dim_eval:n {\paperwidth-25mm},
  height=32mm,
  contents={\makebox[\layerwidth][r]{\_@@_insert_logo:n {\layerheight}}}
]{rwth-logo}
\AddLayersToPageStyle{scrheadings}{rwth-logo}
%    \end{macrocode}
% \iffalse
%</package>
% \fi
% \end{implementation}
\endinput

%</driver>
% \fi
% \begin{documentation}
% \section{Color setup – \pkg{rwth-colors.sty}}
% The colors setup is done using the \pkg{xcolor} package.
% All mechanims described in the corresponding documentation may be used.
%
% All design specific colors are defined according to the guideline \cite{rwth-ci-guideline}.
% The predefined colors are listed below:
%
% \begingroup
% \raggedright
% \ExplSyntaxOn
% \renewcommand*{\definecolor}[3]{
%   \if_mode_vertical:
%   \bool_set_true:N \l_tmpa_bool
%   \tl_range:nnn {#1} {1} {-5}\hfill
%   \tl_range:nnn {#1} {-4} {-1}\space
%   \else:
%   \quad\tl_range:nnn {#1} {-3} {-1}\space
%   \fi:
%   \textcolor{#1}{\rule{1cm}{\ht\strutbox}}
% }
% \begingroup\cctab_select:N \c_document_cctab \input{rwthcolors.def}\endgroup
% \ExplSyntaxOff
% \endgroup
%
% \bigskip
% \noindent As one might realize the shadings are defined directly.
% Those should be used to stick to the CI guideline \cite{rwth-ci-guideline} instead of using the color mixing provided by \pkg{xcolor}.
% \end{documentation}
% \begin{implementation}
% \iffalse
%<*declareoptions>
% \fi
%    \begin{macrocode}
\tl_new:N \g_@@_colormode_tl
\keys_define:nn {ptxcd/colors} {
%    \end{macrocode}
% \iffalse
%</declareoptions>
%<*options>
% \fi
% \end{implementation}
% \begin{documentation}
% \subsection{Package specific options}
% \DescribeKeyOption*{colormode=cmyk/RGB}{Depends on document type}*
% \label{sec:colormode}
% The colormode can be selected to overwrite the default setting.
% Please be aware, that there is a difference in casing and \code{RGB} has to be used uppercase and \code{cmyk} lowercase.
% Internally the values between uppercase and lowercase model names would differ.
% \end{documentation}
% \begin{implementation}
% \begin{optionenv}{colormode}
%    \begin{macrocode}
  colormode .choices:nn = {cmyk,RGB}{
      \@ifpackageloaded{xcolor}{
        \selectcolormodel{\l_keys_value_tl}
      }{
        \exp_args:No \PassOptionsToPackage{\l_keys_value_tl}{xcolor}
        \tl_gset:Ne \g_@@_colormode_tl {\str_lowercase:f {\l_keys_value_tl}}
      }
    },
  colormode .usage:n = load,
  RGB .meta:n = {colormode=RGB},
  cmyk .meta:n = {colormode=cmyk},
%    \end{macrocode}
% \end{optionenv}
%    \begin{macrocode}
}
\ProcessKeyOptions[ptxcd/colors]
\RequirePackage{xcolor}

\input{rwthcolors.def}
%    \end{macrocode}
% \iffalse
%</options>
%<*package&body>
% \fi
%    \begin{macrocode}

%    \end{macrocode}
% \iffalse
%</package&body>
% \fi
% \end{implementation}
% \Finale
\endinput
